\section{The AutoHeal Architecture}
\label{sec:architecture}

AutoHeal is a six-stage closed-loop pipeline executed per maintenance
window. All stage parameters are pre-registered in
\texttt{paper10/design/PAPER10\_PROTOCOL.md} and frozen before any
evaluation-seed AutoHeal result was inspected.

\begin{table}[!t]
  \caption{AutoHeal pipeline: six stages with pre-registered parameters.}
  \label{tab:architecture}
  \centering
  \small
  \setlength{\tabcolsep}{4pt}
  \begin{tabular}{@{}p{1.5cm}p{2.0cm}p{4.0cm}@{}}
    \toprule
    \textbf{Stage} & \textbf{Component} & \textbf{Pre-reg parameters} \\
    \midrule
    Detect       & EEHDA scan      & Real CVE corpus lookup \\
    Score        & HygienePrio     & $(\alpha,\beta,\gamma,\delta)=(0.7,0.5,0.1,0.2)$ \\
    Triage       & 3-class rule    & AUTO $\geq 0.80$; REVIEW $[0.50,0.80)$; DEFER $< 0.50$ \\
    Plan         & Top-$K$         & $K \in \{50,100,200\}$ \\
    Act          & Failure modes   & $0.92$ success / $0.05$ rollback / $0.03$ deferred \\
    Verify       & Health check    & Hard-stop at $>10\%$ rollback or cascade $> 5$ \\
    Learn        & Lag-1 calib     & Applied iff $K \leq 100$; held fixed at $K=200$ \\
    \bottomrule
  \end{tabular}
\end{table}


\subsection{Stage 1 — Detect}

Scan the fleet's vulnerability records; identify unpatched
(host, CVE) pairs. CVE attribute lookup uses the real EPSS / KEV
public corpus (\S\ref{sec:dataset}) so the same CVE in two hosts
has the same external-feed attributes.

\subsection{Stage 2 — Score}

Apply HygienePrio~\cite{paper4hygieneprio} with the calibrated
Paper~4 weights $(\alpha, \beta, \gamma, \delta) = (0.7, 0.5, 0.1,
0.2)$. At capacity $K \leq 100$, online calibration follows
Paper~7's lag-1 recipe~\cite{paper7online}; at $K = 200$, weights
are held fixed per Paper~7's K=200 hazard finding. The scorer
output is normalised to $[0, 1]$ to feed the triage thresholds.

\subsection{Stage 3 — Triage}

Each (host, CVE) pair is classified into one of three buckets via
the pre-registered rule:

\begin{itemize}
\item \textbf{AUTO} (auto-remediate): score $\geq 0.80$ AND
host criticality $\neq$ CRITICAL AND patch-test-suite present AND
no known blocking config AND not during business-hours window.
\item \textbf{REVIEW} (human-in-loop): score in $[0.50, 0.80)$ OR
host criticality = CRITICAL OR business-hours.
\item \textbf{DEFER}: score $< 0.50$ OR known incompatible patch.
\end{itemize}

The two thresholds (0.80, 0.50) and the CRITICAL gate are the
pre-registered configuration. They are deliberately conservative:
AutoHeal only auto-acts on high-confidence, low-risk pairs.

\subsection{Stage 4 — Plan}

From the AUTO bucket, select the top-$K$ pairs by HygienePrio
score (where $K$ is the per-window capacity). This is the
Paper~6~\cite{paper6capacity} capacity-constrained selection.

\subsection{Stage 5 — Act}

Each scheduled pair undergoes a simulated patch action with a
pre-registered failure-mode distribution drawn from public
sysadmin literature on production patch operations:
\begin{itemize}
\item SUCCESS: $92\%$
\item ROLLBACK (post-patch health-check fail): $5\%$
\item DEFERRED (patch blocked at install time): $3\%$
\end{itemize}

The numbers are fixed before evaluation. Real-world deployments
would condition on package, OS, and patch type; AutoHeal's
synthetic evaluation assumes the pre-registered distribution to
keep the threat model identifiable.

\subsection{Stage 6 — Verify + Rollback}

Each acted pair runs a post-action health check. On failure, the
patch is reverted and the pair is re-classified as REVIEW for the
next window. The framework tracks per-window rollback rate; if
$> 10\%$ the framework halts (hard-stop safety bound) and
escalates to human review. A cascading-failure detector
additionally fires if more than 5 rollbacks share a common
upstream patch ancestor (approximated by CVE id prefix).

\subsection{Stage 7 — Learn}

Successful and rolled-back outcomes feed the Paper~7 lag-1
calibration update (when $K \leq 100$). At $K = 200$, weights are
held fixed per Paper~7's K=200 hazard finding, and the
learn stage records data for offline rebaselining only.

\subsection{Pre-Registered Safety Bounds (Hard Stops)}

\begin{itemize}
\item Rollback rate $> 10\%$ at any window $\Rightarrow$ halt;
fall back to human-in-loop for remaining windows.
\item Cascading-failure detection ($> 5$ rollbacks share a common
patch ancestor) $\Rightarrow$ halt; escalate to human.
\item Either bound triggering at any cell-window pair is reported
in the frozen results.
\end{itemize}
